\section{Сборка и проверка}
Проверка интегрированной версии выполнена на Windows: Visual Studio 18 2026, MSVC 19.51.36260.0, Windows SDK 10.0.26100.0, конфигурация Release. Из корня \texttt{Operation-Systems} выполнены команды:
\begin{verbatim}
cmake -S . -B build-lab5
cmake --build build-lab5 --config Release
python lab5/tests/test_lab5.py build-lab5/bin/lab5
\end{verbatim}
Созданы \texttt{lab5\_static.exe}, \texttt{lab5\_dynamic.exe}, \texttt{lab5\_impl1.dll}, \texttt{lab5\_impl2.dll} в \texttt{build-lab5/bin/lab5}. Общая сборка также пересобрала ЛР №2 и №4. Их функциональные тесты не запускались; ЛР №3 пропущена CMake, поскольку требует Linux/WSL.

\subsection{Автоматические проверки}
Тестовый скрипт проверяет пять групп сценариев:
\begin{enumerate}
\item Производную разностью вперёд и площадь прямоугольника в программе №1.
\item Переключение обеих функций на реализацию 2 и обратно: для сторон 3 и 4 результат меняется с 12 на 6 и обратно.
\item Нулевой шаг, отрицательную сторону, \texttt{nan}, переполнение аргумента, нехватку аргументов, лишний аргумент команды \texttt{0}, неизвестную команду и конец ввода без перевода строки.
\item 200 последовательных переключений с загрузкой и освобождением модулей.
\item Перенос файлов в каталог с пробелами и завершение с ошибкой при удалении второй библиотеки перед переключением.
\end{enumerate}
Ниже подключается фактический вывод повторного запуска тестов, сохраняемый при создании отчёта.
\lstinputlisting[language={},numbers=none,caption={Журнал интеграционных проверок Windows}]{lab5/report/inc/tests_windows.txt}

Пять сообщений \texttt{PASS} обозначают группы сценариев, а не пять отдельных вызовов функций. Проверка интегрированной версии в WSL пока не выполнялась. Отдельный тест отсутствующего экспортируемого символа и измерения производительности не проводились.

\section{Анализ способов подключения}
\subsection{Подключение при компоновке}
Вызовы выглядят как обычные вызовы функций; компилятор проверяет аргументы по заголовку, а компоновщик связывает внешние имена с библиотекой. При запуске требуемая DLL должна быть доступна. Изменение реализации при сохранении двоичного контракта не обязательно требует пересборки клиента; изменение самого контракта может потребовать изменения и пересборки~\cite{loadtime}.

\subsection{Загрузка во время исполнения}
Приложение выбирает модуль и момент загрузки самостоятельно. Это позволяет переключать реализации без перезапуска и строить расширяемые приложения. Взамен требуется проверять ошибки загрузки, соответствие имён и время жизни указателей. Отсутствие DLL можно обработать в коде, хотя текущая программа выбирает завершение с диагностикой~\cite{runtime}.

В обоих случаях функции остаются в динамической библиотеке. Отличие заключается в формировании зависимости и управлении загрузкой, а не в противопоставлении статической и динамической библиотеки. Оснований объявлять один способ быстрее другого по выполненным функциональным тестам нет.
